% Lösungen Zone (zusammenbau v0.1)
\einheitenkopf[][Kennst du schon]{Das kennst du schon}

Zuordnung: 1 → der Formelkreis der Einheiten 1 und 2 und die Grundflächen der Einheit 3 · 2 → Einheit 3 und 4 · 3 → Seiten, Höhen und Diagonalen aller Einheiten · 4 → senkrechte Kanten als Höhen erkennen, rechte Winkel prüfen, Einheit 1 bis 3 · 5 → Einheit 3 und 4 · 6 → die Rückrichtungen und Verhältnisse · 7 → der Formelkreis der Einheiten 1 und 2 und die Grundflächen der Einheit 3 · 8 → der Formelkreis der Einheiten 1 und 2 und die Grundflächen der Einheit 3

\erg{1}{a) $20\,\mathrm{cm}^2$ \quad b) $\tfrac{1}{2} \cdot (6{,}5 + 3{,}5) \cdot 4 = 20\,\mathrm{cm}^2$}
\erg{2}{a) $60\,\mathrm{cm}^3$ \quad b) $\pi \cdot 2^2 \cdot 10 \approx 125{,}7\,\mathrm{cm}^3$}
\erg{3}{a) $\overrightarrow{PQ} = (3 | 4 | 0)$ \quad b) $M(1 | 2 | 3{,}5)$}
\erg{4}{a) $0$ – die Vektoren stehen senkrecht \quad b) Skalarprodukt $= 6 - 8 + 2 = 0$, also ja}
\erg{5}{a) $10\,\mathrm{cm}$ \quad b) $\sqrt{6{,}25 - 2{,}25} = 2\,\mathrm{m}$}
\erg{6}{a) $h = 6$ \quad b) $h = 5$}
\erg{7}{a) Lisa nimmt die schräge Seite statt der Höhe. Richtig: $A = \tfrac{1}{2} \cdot 10 \cdot 4 = 20\,\mathrm{cm}^2$}
\erg{8}{a) $\tfrac{1}{2} \cdot 12 \cdot 8 = 48\,\mathrm{cm}^2$}
